\section{Введение: Матрицы Тёплица}
\label{sec:intro}

Вычисление матрично-векторного произведения $\bm{y} = \bm{A}\bm{x}$ для произвольной матрицы $\bm{A} \in \C^{n \times n}$ в общем случае требует $\mathcal{O}(n^2)$ арифметических операций. Однако для структурированных матриц данную сложность можно асимптотически снизить до $\mathcal{O}(n \log n)$. Центральным объектом рассмотрения является класс матриц со специфической структурой диагоналей.

\begin{tcbdefinition}[Матрица Тёплица]
\label{def:toeplitz}
    Матрица $\bm{T} \in \C^{n \times n}$ называется \textit{тёплицевой}, если её элементы постоянны вдоль каждой диагонали. Формально это означает, что элемент матрицы зависит только от разности его индексов: $T_{p,q} = t_{p-q}$. Структура матрицы имеет вид
    \begin{equation}
        \bm{T} =
        \begin{pNiceMatrix}
            t_0     & t_{-1}  & \Cdots & t_{-(n-1)} \\
            t_1     & t_0     & \Ddots & \Vdots     \\
            \Vdots  & \Ddots  & \Ddots & t_{-1}     \\
            t_{n-1} & \Cdots  & t_1    & t_0
        \end{pNiceMatrix}.
        \label{eq:toeplitz}
    \end{equation}
\end{tcbdefinition}

\section{Дискретное преобразование Фурье (ДПФ)}
\label{sec:fft}

Для построения быстрых алгоритмов умножения тёплицевых матриц необходимо ввести аппарат дискретного преобразования Фурье. Во всём дальнейшем тексте предполагается индексация матриц и векторов с нуля.

\begin{tcbdefinition}[Матрица Фурье]
\label{def:fourier_matrix}
    Дискретное преобразование Фурье (ДПФ) вектора $\bm{x} \in \C^n$ задаётся как вектор $\bm{y} \in \C^n$, компоненты которого вычисляются по формуле
    \begin{equation}
        y_p = \sum_{q=0}^{n-1} e^{-\frac{2\pi i}{n} p q} x_q, \quad p = 0, \dots, n-1.
        \label{eq:dft_sum}
    \end{equation}
    В матричной форме преобразование записывается как $\bm{y} = \bm{F}_n \bm{x}$, где $\bm{F}_n$ --- \textit{матрица Фурье} с элементами $(\bm{F}_n)_{p,q} = \omega_n^{pq}$, а $\omega_n = e^{-\frac{2\pi i}{n}}$.
\end{tcbdefinition}

\begin{tcblemma}[Свойства матрицы Фурье]
\label{lem:fourier_props}
    Для матрицы $\bm{F}_n$ выполняются следующие свойства:
    \begin{itemize}
        \item Матрица симметрична: $\bm{F}_n = \bm{F}_n\T$, но не эрмитова (при сопряжении знак в показателе экспоненты меняется).
        \item Выполняется тождество $\bm{F}_n\H \bm{F}_n = n \bm{I}$, следовательно, матрица $\frac{1}{\sqrt{n}} \bm{F}_n$ является унитарной.
    \end{itemize}
\end{tcblemma}

\subsection{Алгоритм быстрого преобразования Фурье (БПФ)}
\label{sec:fft_algo}

При условии, что $n = 2^k, k \in \N$, умножение на матрицу Фурье может быть выполнено за $\mathcal{O}(n \log n)$ операций. Идея заключается в блочном представлении $\bm{F}_{2n}$ через матрицы в два раза меньшего размера $\bm{F}_n$. 

Введём матрицу перестановок $\bm{P}_{2n}$, которая группирует чётные строки исходной матрицы сверху, а нечётные --- снизу. Прямая алгебраическая манипуляция показывает, что
\begin{equation}
    \bm{P}_{2n} \bm{F}_{2n} =
    \begin{pNiceMatrix}
        \bm{F}_n          & \bm{F}_n \\
        \bm{F}_n \bm{W}_n & -\bm{F}_n \bm{W}_n
    \end{pNiceMatrix}
    =
    \begin{pNiceMatrix}
        \bm{F}_n & \bm{0} \\
        \bm{0}   & \bm{F}_n
    \end{pNiceMatrix}
    \begin{pNiceMatrix}
        \bm{I}_n & \bm{I}_n \\
        \bm{W}_n & -\bm{W}_n
    \end{pNiceMatrix},
    \label{eq:fft_block}
\end{equation}
где $\bm{W}_n = \diag(1, \omega_{2n}, \omega_{2n}^2, \dots, \omega_{2n}^{n-1})$.

\begin{tcbalgorithmbox}[Рекурсивное быстрое преобразование Фурье (Radix-2)]
\label{alg:fft}
    \textbf{Вход:} вектор $\bm{x} \in \C^{2n}$, где $2n = 2^k$.\\
    \textbf{Выход:} преобразованный вектор $\bm{y} = \bm{F}_{2n} \bm{x}$.
    \begin{enumerate}
        \item Разделить вектор $\bm{x}$ на две равные части: $\bm{x}_{\text{top}}$ (первые $n$ элементов) и $\bm{x}_{\text{bot}}$ (последние $n$ элементов).
        \item Вычислить базовые суммы и разности (выполнение умножения на правую матрицу из~\cref{eq:fft_block}):
        \begin{enumerate}
            \item $\bm{u} = \bm{x}_{\text{top}} + \bm{x}_{\text{bot}}$.
            \item $\bm{v} = \bm{W}_n (\bm{x}_{\text{top}} - \bm{x}_{\text{bot}})$, где умножение на диагональную матрицу $\bm{W}_n$ выполняется поэлементно.
        \end{enumerate}
        \item Выполнить рекурсивный вызов алгоритма для векторов половинного размера:
              $\bm{y}_{\text{even}} = \bm{F}_n \bm{u}$ и $\bm{y}_{\text{odd}} = \bm{F}_n \bm{v}$.
        \item Переставить элементы результата согласно матрице $\bm{P}_{2n}\T$, поместив $\bm{y}_{\text{even}}$ на позиции чётных индексов, а $\bm{y}_{\text{odd}}$ --- на позиции нечётных.
    \end{enumerate}
\end{tcbalgorithmbox}

\begin{tcbinsight}[Двумерное преобразование Фурье]
    Для двумерного массива $\bm{X} \in \C^{n_1 \times n_2}$ ДПФ определяется как $\bm{Y} = \bm{F}_{n_1} \bm{X} \bm{F}_{n_2}\T$. В терминах векторизации матриц данная операция описывается через кронекерово произведение:
    \[
        \operatorname{vec}(\bm{Y}) = (\bm{F}_{n_2} \otimes \bm{F}_{n_1}) \operatorname{vec}(\bm{X}).
    \]
\end{tcbinsight}

\section{Циркулянты и дискретная теорема о свёртке}
\label{sec:circulant}

\begin{tcbdefinition}[Циркулянт]
\label{def:circulant}
    Тёплицева матрица $\bm{C} \in \C^{n \times n}$ называется \textit{циркулянтом}, если каждый её последующий столбец является циклическим сдвигом предыдущего вниз:
    \begin{equation}
        \bm{C} =
        \begin{pNiceMatrix}
            c_0     & c_{n-1} & \Cdots & c_1 \\
            c_1     & c_0     & \Cdots & c_2 \\
            \Vdots  & \Vdots  & \Ddots & \Vdots \\
            c_{n-1} & c_{n-2} & \Cdots & c_0
        \end{pNiceMatrix}.
        \label{eq:circulant}
    \end{equation}
\end{tcbdefinition}

Умножение вектора на циркулянт $\bm{y} = \bm{C}\bm{x}$ порождает операцию, известную как \textit{периодическая свёртка}. Элементы вектора результата задаются формулой $y_p = \sum_{q=0}^{n-1} c_{(p-q) \bmod n} \, x_q$.

\begin{tcbtheorem}[Спектральное разложение циркулянта]
\label{thm:circulant_eig}
    Любой циркулянт $\bm{C}$ с первым столбцом $\bm{c}$ диагонализуется матрицей Фурье $\bm{F}_n$:
    \begin{equation}
        \bm{C} = \bm{F}_n^{-1} \diag(\bm{F}_n \bm{c}) \bm{F}_n.
        \label{eq:circulant_eig}
    \end{equation}
    Иными словами, собственными значениями циркулянта являются компоненты вектора $\bm{F}_n \bm{c}$, а матрица собственных векторов универсальна и совпадает с $\bm{F}_n^{-1}$.
\end{tcbtheorem}

\begin{tcbproof}[теоремы~\ref{thm:circulant_eig}]
    Пусть $\eta$ --- комплексное число такое, что $\eta^n = 1$. Рассмотрим полином $\lambda(\eta) = c_0 + \eta c_1 + \dots + \eta^{n-1} c_{n-1}$.
    Вычислим произведение вектора-строки $(1, \eta, \eta^2, \dots, \eta^{n-1})$ на циркулянт $\bm{C}$ слева:
    \begin{equation*}
        (1, \eta, \dots, \eta^{n-1}) \bm{C} = \lambda(\eta) (1, \eta, \dots, \eta^{n-1}).
    \end{equation*}
    Полагая $\eta_k = \omega_n^k$ для $k = 0, \dots, n-1$, мы получаем $n$ левых собственных векторов матрицы $\bm{C}$, которые в точности образуют строки матрицы Фурье $\bm{F}_n$. Отсюда следует матричное равенство $\bm{F}_n \bm{C} = \bm{\Lambda} \bm{F}_n$, где $\bm{\Lambda} = \diag(\lambda(\eta_0), \dots, \lambda(\eta_{n-1}))$. Подстановка определения $\lambda(\eta_k)$ показывает, что вектор диагональных элементов есть в точности $\bm{F}_n \bm{c}$.
\end{tcbproof}

\begin{tcbcorollary}[Дискретная теорема о свёртке]
\label{cor:discrete_conv}
    Для вычисления периодической свёртки (умножения циркулянта $\bm{C}$ на вектор $\bm{x}$) справедлива формула:
    \begin{equation}
        \bm{C}\bm{x} = \bm{F}_n^{-1} \bigl( (\bm{F}_n \bm{c}) \circ (\bm{F}_n \bm{x}) \bigr),
    \end{equation}
    где $\circ$ обозначает поэлементное произведение (произведение Адамара).
\end{tcbcorollary}

В силу \cref{cor:discrete_conv}, умножение циркулянта на вектор требует вычисления двух прямых и одного обратного БПФ, а также одной операции поэлементного перемножения векторов. Следовательно, общая сложность операции составляет $\mathcal{O}(n \log n)$.

\section{Быстрое умножение матрицы Тёплица}
\label{sec:toeplitz_fast}

Произвольная тёплицева матрица не обладает цикличностью, однако её можно искусственно погрузить в циркулянт большего размера, чтобы воспользоваться быстрым алгоритмом на базе БПФ.

\begin{tcblemma}[Вложение Тёплица в циркулянт]
\label{lem:embedding}
    Любую тёплицеву матрицу $\bm{T} \in \C^{n \times n}$ можно вложить в качестве левого верхнего блока в циркулянт $\bm{C}$ размера $N \times N$, где $N \geq 2n-1$.
\end{tcblemma}

\begin{tcbalgorithmbox}[Умножение тёплицевой матрицы на вектор]
\label{alg:toeplitz_mult}
    \textbf{Вход:} тёплицева матрица $\bm{T} \in \C^{n \times n}$ и вектор $\bm{x} \in \C^n$.\\
    \textbf{Выход:} вектор $\bm{y} = \bm{T}\bm{x}$.
    \begin{enumerate}
        \item Выбрать размер циркулянта $N = 2^q \geq 2n-1$, чтобы $N$ являлось степенью двойки для оптимальной работы БПФ.
        \item Построить первый столбец $\bm{c} \in \C^N$ расширенного циркулянта $\bm{C}$:
        \[
            \bm{c} = (t_0, t_1, \dots, t_{n-1}, 0, \dots, 0, t_{-(n-1)}, \dots, t_{-1})\T.
        \]
        \item Построить расширенный вектор $\bm{\tilde{x}} \in \C^N$, дописав нули:
        \[
            \bm{\tilde{x}} = (x_0, x_1, \dots, x_{n-1}, 0, \dots, 0)\T.
        \]
        \item Вычислить произведение $\bm{\tilde{y}} = \bm{C}\bm{\tilde{x}}$ с помощью дискретной теоремы о свёртке (\cref{cor:discrete_conv}):
        \[
            \bm{\tilde{y}} = \bm{F}_N^{-1} \bigl( (\bm{F}_N \bm{c}) \circ (\bm{F}_N \bm{\tilde{x}}) \bigr).
        \]
        \item Искомый вектор $\bm{y}$ совпадает с первыми $n$ компонентами вектора $\bm{\tilde{y}}$.
    \end{enumerate}
\end{tcbalgorithmbox}

\section{Вычисление БПФ произвольного размера}
\label{sec:bluestein}

Алгоритм \cref{alg:fft} жёстко требует, чтобы размер преобразования был степенью двойки. Для произвольного размера $n$ применяется сведение задачи ДПФ к умножению на тёплицеву матрицу на основе алгебраического тождества
\begin{equation}
    pq = \frac{p^2 + q^2 - (p-q)^2}{2}.
    \label{eq:bluestein_id}
\end{equation}

Подстановка \cref{eq:bluestein_id} в базовую формулу ДПФ из \cref{def:fourier_matrix} приводит к следующей факторизации суммы:
\begin{align}
    y_p &= \sum_{q=0}^{n-1} e^{-\frac{2\pi i}{n} pq} x_q \notag \\
    &= \omega_n^{\frac{p^2}{2}} \sum_{q=0}^{n-1} \omega_n^{-\frac{(p-q)^2}{2}} \left( \omega_n^{\frac{q^2}{2}} x_q \right).
    \label{eq:bluestein_sum}
\end{align}

\begin{tcbinsight}[Трюк Блюштейна]
    Равенство~\cref{eq:bluestein_sum} описывает разбиение умножения $\bm{F}_n \bm{x}$ на три этапа:
    \begin{enumerate}
        \item Поэлементное умножение: вектор $\bm{x}$ слева умножается на диагональную матрицу с элементами $\omega_n^{q^2/2}$.
        \item Умножение на тёплицеву матрицу: полученный вектор умножается на матрицу $\bm{T}$, элементы которой $T_{p,q} = \omega_n^{-(p-q)^2/2}$.
        \item Поэлементное умножение: результат умножается на диагональную матрицу с элементами $\omega_n^{p^2/2}$.
    \end{enumerate}
\end{tcbinsight}

Умножение на полученную тёплицеву матрицу осуществляется с помощью \cref{alg:toeplitz_mult} с вложением в циркулянт размера $N = 2^k \geq 2n-1$. Таким образом, вычисление ДПФ произвольного размера алгоритмически сводится к операциям БПФ в размере степени двойки, сохраняя асимптотическую сложность $\mathcal{O}(N \log N) = \mathcal{O}(n \log n)$.
